\section{Introduction}\label{sec:intro}

A chemical system with several stable compositions can, in principle, store information in the composition it occupies. This idea underlies compositional models of prebiotic heredity \citep{segre2000,vasas2012}, analyses of inheritance in compartmentalized autocatalytic networks under growth and division \citep{matsubara2026,singh2023}, and the question of how many heritable states a chemistry with a given number of species can support \citep{virgo2015}. Bistable and oscillatory networks of small organic molecules have been realized experimentally in flow reactors \citep{semenov2016}. None of this establishes that a \emph{finite} population of molecules copies a stored state reliably, and the gap is not a technicality. Growth changes concentrations and may consume the very species that supports a state. Division allocates molecules at random, and the two daughters are complementary rather than independent. Any recovery after birth takes time and can itself fail. A stable composition is therefore only the starting point; what has to be controlled is the complete cycle.

This paper asks a quantitative version of the question. Under a fully specified stochastic growth-and-division protocol, can both daughters recover the parent's chemical label and become admissible parents themselves, with an explicit probability and an explicit molecular budget? And when more labels must be stored, or more divisions must succeed, how does the required molecular redundancy grow?

\paragraph{Copying as a complete cycle.}
Throughout, a label $\ell$ is represented by an \emph{admitted region} $B_\ell$ of newborn or restart states, chosen so that the copying guarantee is uniform over $B_\ell$. One cycle succeeds when an admitted parent reaches the specified division event within the allowed time and resources, its molecules are distributed to two daughters by the protocol's actual partition law, and both daughters satisfy the same label's admission condition at the protocol's inspection time. Success is deliberately stronger than a decoder assigning the correct label: a daughter can pass a threshold test while lying outside the set on which the next cycle's theorem applies. The reward for this stronger definition is that repeated copying follows by conditioning alone, with no independence assumption between siblings or generations (Section~\ref{sec:events}).

\paragraph{Three results.}
We prove three separate results and keep their scopes distinct; Table~\ref{tab:three} summarizes them.
\begin{enumerate}[label=(\Alph*)]
\item \emph{A small constructed source} (Section~\ref{sec:constructed}). Two interacting reversible two-species modules, coupled through shared consumptive growth and cross-catalysis, store four concentration-state labels with at most $874$ resident molecules in each newborn module. Under an ideal external controller that prescribes the volume and holds two food and two waste pools at fixed activities, every admitted parent produces two correctly returned daughters at its first division within $20$ seconds with probability at least $0.9901$, uniformly over the interaction strength $\epsilon\in[0,0.1]$ and the growth-reversal rate $\rho\in[0,2\times10^{-7}]$, and in particular on a strictly positive rectangle. The proof is a finite-state generator certificate checked in exact integer arithmetic, together with a sharp stopped-counter estimate that makes the complete-cycle reservoir capacity per pool about $0.166$ pmol.
\item \emph{Molecular redundancy for many modules} (Section~\ref{sec:redundancy}). For $k$ interacting modules with $d$ resident species each, sharing consumptive growth and exchanging a growth species through a graph with bounded total incident weight, a local quadratic recovery inequality and bounded reaction activity give a one-generation failure probability at most $kC(\gamma)e^{-cN}$, where $N$ is the birth membrane count per module and the exponent $c$ and interaction allowance do not depend on $k$. For an explicit four-species module the sufficient copy number for a designated lineage of $G$ divisions is $O(\log(kG/\eta))$ at target error $\eta$, and $O(G+\log(k/\eta))$ for a complete binary family. A partition obstruction gives matching necessary orders for immediate return of both daughters, so these orders are sharp for that copying criterion. The constants are astronomically conservative, and the sufficient bound rescales physical rates with $k$; we state exactly which costs the asymptotic statement hides.
\item \emph{A nominal published chemistry} (Section~\ref{sec:semenov}). The eight-species kinetic model reconstructed from Semenov et al.\ \citep{semenov2016}, at its nominal rate constants, supports a prescribed protocol of fair splitting, refilling with fresh feed and $500$ seconds of recovery in a continuously stirred reactor of $39.853\,\micro$L: both daughters return to the same rational return region with probability at least $124/125=0.992$. The certified regions are narrow, the recovery is imposed, and deterministic sensitivity diagnostics show why the theorem must not be read as an experimentally validated uncertainty range.
\end{enumerate}
The conceptual link between the three is reliable return of a label under specified operations. They do not combine into one chemistry that is simultaneously small, experimentally identified and autonomous: the $874$-molecule bound belongs to the constructed source, the logarithmic orders belong to the redundancy theorem's own chemistry, and the Semenov result uses roughly $10^{18}$ active molecules under an external protocol.

\begin{table}[t]
\centering\small
\caption{The three results and their scopes. ``Return'' always means that both daughters lie in the admitted region of the same label; the inspection time differs between protocols.}\label{tab:three}
\begin{tabular}{@{}>{\raggedright\arraybackslash}p{2.6cm}>{\raggedright\arraybackslash}p{3.9cm}>{\raggedright\arraybackslash}p{3.9cm}>{\raggedright\arraybackslash}p{3.9cm}@{}}
\toprule
 & (A) Constructed source & (B) Redundancy theorem & (C) Nominal Semenov model\\
\midrule
Chemistry & two two-species modules, hypothetical reversible rate table, shared growth marker, cross-catalysis & $k$ modules of $d$ species; explicit four-species and two-species instances; growth-species exchange graph & eight tracked species of Semenov et al.\ \citep{semenov2016} in a flow reactor at nominal rates\\
Stored quantity & four concentration-state words $\sigma\in\{0,1\}^2$ & words $\sigma\in\{0,1\}^k$ with a relative-composition readout & two labels, thiol concentration state (relative readout also valid)\\
Resource measure & $\le874$ residents per newborn module; birth marker $N=53$; pools reset by the controller & birth membrane count $N$ per module; $N\ge1.4\times10^{20}$ floor for the four-species instance & $\Omega=2.4\times10^{19}$ molecules per molar, $V=39.853\,\micro$L; feed quota $1.2\times10^{19}$ per daughter\\
Success event & first division by $20$ s, no monitor triggered, both daughters in the boxes at birth & certified generation, or immediate first division, with both daughters in the birth regions & both daughters in the return ellipsoid after split, refill and $500$ s recovery\\
Guarantee & $\ge0.9901$ per cycle, uniform in $\epsilon\in[0,0.1]$, $\rho\in[0,2\times10^{-7}]$ & failure $\le kC(\gamma)e^{-cN}$; $N=\Theta(\log(kG/\eta))$ lineage, $\Theta(G+\log(k/\eta))$ family & $\ge0.992$ per cycle, uniform over the return regions\\
Main limitations & ideal instantaneous pool and volume control; hypothetical compounds & conservative constants; rates rescaled with $k$; deadline $O(k/\gamma)$; necessity only for immediate return & nominal parameters; narrow regions; imposed recovery; known model/data discrepancy\\
Verification & Lean~4, exact integer tables & Lean~4 for the certified kernel and the necessity bound; one conventional transfer lemma & Lean~4, exact rational polynomial certificates\\
\bottomrule
\end{tabular}
\end{table}

\paragraph{What is not claimed.}
None of the results establishes autonomous compartment growth, membrane geometry, selection, differential reproduction, or open-ended evolution; the guarantees concern copying under the stated external operations. The Semenov result is a prediction of a specified nominal model, not an experimental demonstration. Neither result claims to introduce chemical heredity or bistability as topics. We also do not claim that exponential residence at large volume is new; it has general antecedents in large-deviation theory for reaction networks \citep{agazzi2018} and in Lyapunov methods for stochastic reaction networks \citep{gupta2014,kuntz2019}. The contributions are the explicit, uniform finite-copy bounds attached to literal stochastic protocols, the uniformity in the number of modules, and the fact that every quantitative constant is the output of an exactly checked certificate rather than an estimate from simulated trajectories.

\paragraph{Methods.}
All processes are continuous-time Markov jump processes on integer counts with mass-action propensities \citep{gillespie1977,kurtz1972,andersonkurtz2015}. Probabilities of one cycle are bounded by comparison functions on a finite stopped model: a subsolution $v$ of the backward inequality and a nonnegative decaying witness $w$ turn a generator inequality into a deadline bound (Lemma~\ref{lem:deadline}). The comparison functions are proposed numerically and then checked in exact integer or rational arithmetic; for the Semenov model they are time-dependent rational quadratic forms on $38$ polynomial time pieces. Division is handled through the actual complementary binomial allocation, either as an exact terminal payoff or through union bounds, never through a product of sibling probabilities. The passage from finite stopped models to the actual count process uses nonexplosion and a pathwise coupling up to the first monitored event. The finite-model theorems and the certified-generation kernel of the redundancy theorem are checked in Lean~4 with Mathlib \citep{demoura2021,mathlib2020}; Appendix~\ref{app:verification} lists the declarations, the receipts and the exact boundary of what is machine-checked.

\paragraph{Organization.}
Section~\ref{sec:events} defines the copying cycle, offspring kernels and the repeated-copying lemma, and distinguishes a designated lineage from a complete family. Section~\ref{sec:constructed} states and proves the constructed-source theorem. Section~\ref{sec:redundancy} states the general transmission principle, its four-species instance, the immediate-copying bounds and the resulting molecular-redundancy orders. Section~\ref{sec:semenov} treats the nominal Semenov protocol. Section~\ref{sec:discussion} discusses what the results establish and what they leave open. The appendices give the verification map, the four-species source certificates and the Semenov certificate constants.
